\section{Related Work}

\paragraph{Diffusion models and latent diffusion.}
Denoising diffusion probabilistic models learn to reverse a Gaussian noising
process by predicting noise targets \citep{ho2020denoising}. Latent diffusion
models move this process into a learned latent space and decode the final latent
through a VAE-style decoder \citep{rombach2022latentdiffusion}. \method{}
instead stays in the pixel-space denoising algebra and factorizes the dense
noise prediction itself. The \method{} token is not a final image latent; it is a
negative-noise component that directly contributes to the pixel-space
denoising update.

\paragraph{Trajectory ordering and latent-pixel interaction.}
Latent Forcing studies how changing the schedule between latent and pixel
variables changes the information revealed by the diffusion trajectory
\citep{baade2026latentforcing}. \method{} is inspired by this view of trajectory
ordering, but tokenizes the noise prediction inside the reverse trajectory
rather than introducing a separate latent-pixel trajectory.

\paragraph{Ordered visual tokenizers.}
Spectral Image Tokenizer, FlexTok, Matryoshka Representation Learning, M3, and
related nested-token approaches organize image or multimodal representations so
that prefixes or subsets can be useful
\citep{esteves2025spectral,bachmann2025flextok,kusupati2024matryoshka,cai2024matryoshkamultimodal}.
These methods operate on image representations or visual-token budgets.
\method{} operates on dense diffusion noise prediction. Its prefixes are valid
pixel-space denoising updates, not partial decodes from a general image
tokenizer.

\paragraph{Diffusion-autoregressive visual tokens.}
Selftok and D-AR are the closest conceptual neighbors because they connect
diffusion processes and autoregressive tokens \citep{wang2025selftok,gao2025dar}.
\method{} differs in the represented object and in the anti-degeneration
mechanism. \method{} tokens are continuous negative-noise components,
and the synthesis operator is deliberately weak and condition-free. The main
claim is not that diffusion steps are visual tokens, but that dense noise
prediction can be decomposed into ordered components whose prefixes are directly
usable in the diffusion denoising formula.
